\secnumbersection{CONCLUSIONES}

\javi{Profe! no se como citar los capítulos, intente colocando label al lado pero me cita la ultima seccion del capitulo anterior}
\eli{Me parece que es por cómo está redefinido el entorno. No rabees más con eso y déjalo fijo.}


Esta memoria abordó el desarrollo y la evaluación de un modelo basado en Aprendizaje por Refuerzo (RL) para la resolución del Problema de Ruteo de Vehículos Periódico (PVRP), planteado como una alternativa data-driven a las metaheurísticas clásicas que tradicionalmente han dominado la resolución de esta familia de problemas. El trabajo se estructuró en torno a la reformulación del PVRP como un Proceso de Decisión de Markov (MDP), permitiendo que un agente construya soluciones de manera secuencial mediante la interacción con un entorno de simulación, en lugar de depender de reglas de decisión diseñadas manualmente. La propuesta contempló cuatro componentes de diseño principales: 

\begin{itemize}
    \item La representación del estado como un vector de dimensión fija que integra información espacial y temporal.
    
    \item Un espacio de acciones discreto que incluye explícitamente la decisión de cerrar la ruta actual, otorgando al agente control directo sobre la estructura de la solución.
    
    \item Una función de recompensa con penalización proporcional al número de clientes no atendidos.
    
    \item Un mecanismo de enmascaramiento de acciones que garantiza la factibilidad local en cada paso de decisión, reduciendo el espacio de exploración durante el entrenamiento.
\end{itemize}
 
Sobre esta formulación se implementó un agente basado en el algoritmo Maskable PPO, entrenado y evaluado sobre nueve instancias del conjunto NEO~\cite{NEO2013} agrupadas en tres familias estructurales, con evaluación estadísticamente rigurosa mediante análisis multi-semilla sobre siete de esas instancias (cinco semillas cada una, para un total de treinta y cinco entrenamientos independientes), contrastada con dos métodos de referencia: una heurística voraz (Greedy) y la metaheurística Variable Neighborhood Search (VNS), reconocida en la literatura como uno de los métodos más efectivos para el PVRP~\cite{HEMMELMAYR2009791}. Adicionalmente, se evaluó la transferencia directa entre instancias y el entrenamiento conjunto sobre múltiples instancias, ejes que en conjunto permiten caracterizar no solo la calidad de las soluciones obtenidas, sino también la robustez estadística del método y su capacidad de generalización más allá de la instancia específica de entrenamiento.

Para evaluar el éxito de esta memoria, resulta fundamental contrastar los resultados obtenidos con los compromisos metodológicos iniciales. A lo largo del desarrollo, se cumplió con el objetivo general planteado en la propuesta original: desarrollar y evaluar un modelo basado en Aprendizaje por Refuerzo para la resolución del PVRP, con el propósito de determinar políticas de ruteo eficientes y adaptativas. Asimismo, se abordaron con éxito los cinco objetivos específicos comprometidos, logrando en varios casos superar ampliamente el alcance mínimo estipulado.

A continuación, se detalla el cumplimiento metodológico de cada uno de estos hitos:

\begin{itemize}
    \item \textbf{Revisión del estado del arte:} Este primer objetivo se realizó en el Capítulo~1 mediante una investigación bibliográfica profunda. Esta revisión abarcó desde los métodos exactos y heurísticos clásicos hasta los enfoques híbridos más recientes de Aprendizaje por Refuerzo aplicados al ruteo, lo que permitió situar la propuesta de esta memoria dentro de una línea de investigación científica altamente activa y en constante desarrollo.
    
    \item \textbf{Creación del entorno de simulación:} El segundo objetivo demandaba modelar adecuadamente las características y restricciones del PVRP. Para esto se materializó esta meta construyendo un entorno completamente funcional que reproduce de manera fiel el Proceso de Decisión de Markov descrito en el Capítulo~2. Con el fin de garantizar la robustez del software, se validó mediante $142$ pruebas automatizadas. Estas pruebas verifican estrictamente tanto las invariantes locales de factibilidad como la coherencia absoluta entre el costo acumulado y el costo oficial reportado por el conjunto NEO.
    
    \item \textbf{Desarrollo del agente funcional:} El tercer objetivo exigía un agente capaz de interactuar con el entorno y completar un ciclo de entrenamiento. Se realizó este hito implementando un agente basado en el algoritmo Maskable PPO. Calibrando este modelo de forma sistemática sobre la instancia base y, posteriormente, se recalibró para abordar escalas mayores. En esta etapa se superó el alcance mínimo estipulado en la propuesta original, ya que se diseñaron e incluyeron experimentos adicionales de transferencia y entrenamiento multi-instancia que no formaban parte de los compromisos iniciales.
    
    \item \textbf{Evaluación del modelo entrenado:} El cuarto objetivo apuntaba a medir el gap porcentual del modelo respecto al mejor resultado conocido (BKS). Esta tarea se abordó evaluando el agente sobre nueve instancias del conjunto NEO, cubriendo tres escalas de tamaño y tres familias estructurales distintas. Para otorgar un respaldo estadístico riguroso a los resultados, se utilizó una metodología multi-semilla, lo que permitió distinguir claramente el comportamiento reproducible del método frente a la variabilidad aleatoria propia de una corrida individual.
    
    \item \textbf{Comparación con baselines de la literatura:} Finalmente, el quinto objetivo requería contrastar la propuesta con métodos de resolución ya reportados. Este requisito se satisfizo mediante la implementación y ejecución de una heurística Greedy y de la metaheurística VNS sobre las nueve instancias de prueba. Con esto, se logró contrastar de forma exhaustiva la calidad de las soluciones, los tiempos de ejecución y las tasas de factibilidad en todos los escenarios, yendo mucho más allá de una simple validación sobre la instancia base.
\end{itemize}

Los experimentos realizados permiten exponer un conjunto de hallazgos centrales, todos sustentados estrictamente en la evidencia empírica reportada en el Capítulo~4. En primer lugar, destaca el rendimiento obtenido en las instancias de $50$ clientes pertenecientes a las familias A y C, donde el agente propuesto supera con claridad a la metaheurística VNS. Específicamente, en la instancia $p01$ el modelo de RL logra un gap medio de $+19{,}7\%$ (con una desviación de $\pm 6{,}5\%$ sobre cinco semillas), superando el $+31{,}7\%$ alcanzado por VNS. Esta superioridad se acentúa en la instancia $p03$, donde el gap del agente se sitúa en $+28{,}0\%$, marcando una diferencia notable frente al $+83{,}8\%$ de la metaheurística. Por su parte, la heurística voraz queda relegada a márgenes muy inferiores, registrando $+70{,}2\%$ y $+100{,}2\%$ en dichas instancias. Esta brecha de rendimiento confirma que el aprendizaje de una política basada en la experiencia acumulada aporta un valor analítico sustancial por sobre una regla de construcción miope. En consecuencia, se evidencia el alto valor de un enfoque aprendido para la construcción de rutas en escenarios de tamaño reducido.

Más allá de la calidad geométrica de las soluciones, el resultado más innovador de la evaluación radica en la eficiencia computacional durante la fase de inferencia. El agente entrega rutas completas entre uno y dos órdenes de magnitud más rápido que los métodos iterativos de comparación. Mientras que VNS requiere tiempos de ejecución que oscilan entre décimas de segundo y $38{,}7$ segundos dependiendo de la complejidad de la instancia, el modelo entrenado responde de manera consistente en apenas fracciones de segundo. Esta ventaja operativa trasciende el análisis académico y cobra particular relevancia en contextos operativos reales, donde la rapidez de respuesta condiciona directamente la viabilidad del despliegue tecnológico. En sistemas de planificación diaria en logística de distribución, la capacidad de recalcular rutas de forma instantánea ante cambios de última hora o contingencias en terreno resulta crítica, exigiendo tiempos de respuesta muy inferiores a los que una metaheurística iterativa tradicional puede ofrecer.

El método también escala estructuralmente a instancias de $75$ y $100$ clientes, manteniendo factibilidad en veintinueve de treinta corridas evaluadas (tasa global del $97\%$ en las familias A y C), aunque la calidad relativa al BKS se degrada con el tamaño: en la familia A, el gap medio pasa de $+19{,}7\%$ en $p01$ ($50$ clientes) a $+53{,}9\%$ en $p04$ ($75$ clientes) y $+57{,}4\%$ en $p07$ ($100$ clientes), con un patrón análogo en la familia C. El salto se produce principalmente entre $50$ y $75$ clientes, mientras que entre $75$ y $100$ la degradación se mitiga, sugiriendo una estabilización del método en esas escalas más que un deterioro indefinido a medida que crece la instancia. Este comportamiento es, además, reproducible: la desviación estándar del gap entre semillas se mantiene en el rango $\pm 5\%$ a $\pm 7\%$ en todas las configuraciones evaluadas, independientemente de la escala, lo que indica que los resultados observados no son producto del azar de una ejecución en particular sino una propiedad estable del método bajo la configuración adoptada, y permite descartar interpretaciones triunfalistas o pesimistas basadas en corridas individuales.

Junto a estas fortalezas, la evaluación identificó un límite combinatorio claro en la familia B: el agente no alcanza factibilidad en $p02$, instancia que el VNS resuelve sin dificultad ($+20{,}0\%$, factible; Sección~\ref{sec:limite-familiaB}). Descartada la saturación de capacidad como causa, dado que $p03$ posee saturación mayor y sí se resuelve, el fenómeno se atribuye a la naturaleza secuencial del enfoque constructivo sin retroceso, incapaz de coordinar simultáneamente la asignación multi-frecuencia y el reparto de carga entre múltiples vehículos. Los experimentos de transferencia confirmaron que este límite es robusto a la estrategia de entrenamiento, ya que persiste al evaluar sobre $p02$ agentes entrenados en otras instancias.

Los experimentos de transferencia directa (Sección~\ref{sec:transferencia}) entregaron uno de los hallazgos más matizados de esta investigación. Al evaluar cruzadamente los agentes especializados en $p01$ y $p03$, que comparten la geometría de clientes pero difieren en horizonte temporal y número de vehículos, se observaron dos regímenes opuestos. En el sentido $p01 \rightarrow p03$, la política reprodujo su solución paso a paso de forma ciega al nuevo contexto, produciendo el mismo costo ($617{,}39$) y resultando factible solo por coincidencia estructural entre el presupuesto de rutas de ambas configuraciones, lo que constituye una transferencia contingente y no una adaptación real. En el sentido inverso $p03 \rightarrow p01$, en cambio, la secuencia de decisiones divergió al percibir el cambio en la variable de estado que codifica el vehículo activo, construyendo una partición distinta que resultó factible e incluso mejoró el gap base ($+20{,}1\%$ frente a $+21{,}3\%$), lo que sí evidencia una adaptación genuina.

En conjunto, ambos experimentos demuestran que el agente es capaz de generalizar su comportamiento siempre y cuando la diferencia de contexto sea percibida correctamente por su representación de estado. Asimismo, dejan una lección metodológica de gran relevancia para la evaluación de sistemas de aprendizaje aplicados a la optimización combinatoria: resulta obligatorio inspeccionar el mecanismo secuencial detrás del costo final antes de catalogarlo como un aprendizaje transferible. Finalmente, a estos hallazgos se suma el análisis del entrenamiento conjunto sobre las tres instancias, denominado \texttt{ppo\_multi\_p01\_p02\_p03}. Si bien este enfoque produjo soluciones factibles en las familias A y C, experimentó una degradación de calidad muy significativa respecto a los agentes especializados, arrojando gaps que oscilaron entre el $+58\%$ y el $+94\%$. Esta caída en el rendimiento, equivalente a resultados aproximadamente cuatro veces peores, pone en evidencia la fuerte tensión existente entre la especialización y la generalización, una característica recurrente al aplicar algoritmos de RL en problemas de alta complejidad combinatoria.

Más allá de los resultados numéricos de calidad y generalización, el proceso de desarrollo del agente dejó dos lecciones metodológicas que resulta fundamental registrar explícitamente, pues ilustran la alta sensibilidad del comportamiento de un agente de Aprendizaje por Refuerzo frente a decisiones de diseño aparentemente secundarias. La primera de ellas se refiere a la calibración de hiperparámetros. Al utilizar los valores por defecto del algoritmo, consistentes en un coeficiente de entropía de $0{,}01$ y una red neuronal de $[128, 128]$ neuronas, la política convergía prematuramente hacia soluciones de baja calidad, quedando atrapada en óptimos locales tempranos. Fue necesario aumentar el coeficiente de entropía a $0{,}05$ para incentivar una exploración más prolongada, junto con ampliar la red a $[256, 256]$ neuronas para otorgar una mayor capacidad de representación al procesar el vector de estado del PVRP, logrando así la configuración que finalmente permitió superar a los métodos de referencia clásicos. Esta sensibilidad algorítmica confirma que, en la aplicación de modelos de aprendizaje a problemas de optimización combinatoria, la calibración de hiperparámetros trasciende de ser un detalle de afinación menor para convertirse en una condición estrictamente necesaria en la obtención de resultados competitivos. La segunda lección, de una relevancia aún mayor, surgió durante el desarrollo del entrenamiento multi-instancia y constituye un caso claro de vulneración de la recompensa o reward hacking. La función de recompensa original penalizaba la infactibilidad con un valor fijo, independientemente de la cantidad de clientes que quedaran sin atender, provocando que abandonar a un solo cliente costara computacionalmente lo mismo que abandonar a la totalidad de la demanda. El agente descubrió y explotó rápidamente este atajo, colapsando hacia soluciones triviales compuestas por una sola ruta y un único cliente visitado, lo que generaba un costo aparente mínimo pero producía una tasa de factibilidad que se desplomaba a cero. La corrección metodológica consistió en estructurar una penalización estrictamente proporcional al número de clientes no visitados, ajuste tras el cual la tasa de episodios factibles del entrenamiento multi-instancia ascendió y se estabilizó de forma sostenida, sin registrar nuevos colapsos. Este episodio ilustra un principio general crítico en el diseño de funciones de recompensa: el agente optimiza de manera literal la señal matemática que se le entrega y no la intención subyacente del diseñador, por lo que una penalización que no logra distinguir los grados de severidad frente a la violación de una restricción invita invariablemente al modelo a explotar el camino más barato para evadirla.

La rigurosidad científica exige delimitar con precisión los alcances de esta investigación y reconocer explícitamente las limitaciones del trabajo desarrollado. En términos de alcance empírico, se logró una caracterización estadísticamente robusta del agente sobre nueve instancias del conjunto NEO, abarcando tres escalas de tamaño y tres familias estructurales distintas bajo una implementación completamente reproducible. Esta reproducibilidad se garantiza mediante la disponibilidad de los modelos entrenados, un código de software respaldado por $142$ pruebas automatizadas y la capacidad de ejecutar la totalidad de los experimentos sobre hardware estándar sin requerir aceleración gráfica mediante GPU, lo que facilita la replicación del estudio sin depender de infraestructura computacional especializada. En cuanto a las limitaciones metodológicas, se constata que el agente construye las soluciones de forma estrictamente incremental sin contemplar una fase de mejora posterior. Esta carencia algorítmica se manifiesta visualmente en cruces de aristas que una optimización tradicional de tipo $2$-opt eliminaría con facilidad, lo que explica en gran medida la pérdida de calidad frente a la metaheurística VNS al escalar el tamaño de la red. Además, la arquitectura de la red neuronal opera sobre vectores de estado cuya dimensión está rígidamente ligada al número de clientes, impidiendo la transferencia directa de políticas aprendidas entre instancias de distinto tamaño sin someter al modelo a un reentrenamiento previo. Asimismo, la evaluación empírica se restringió al conjunto NEO estándar, marginando variantes más complejas que incluyen demandas estocásticas, ventanas de tiempo u otras restricciones dinámicas que son habituales en escenarios logísticos reales. Respecto a las limitaciones inherentes al enfoque de resolución propuesto, destaca el hecho de que el agente no logró alcanzar la factibilidad en la única instancia evaluada de la familia B, correspondiente a $p02$. Aunque las instancias $p05$ y $p08$ de esa misma categoría estructural no fueron entrenadas explícitamente debido al elevado costo computacional que demanda el protocolo multi-semilla a esa escala, resulta altamente razonable que el enfoque constructivo puro enfrente dificultades operativas análogas, configurando una limitación estructural del método y no un simple defecto de implementación. De igual forma, en las instancias de $75$ y $100$ clientes, el VNS supera al agente en la calidad de la ruta, evidenciando que la aproximación constructiva no representa la estrategia óptima en dichas escalas si el objetivo es la estricta minimización del costo, aun conservando su ventaja competitiva en el tiempo de inferencia. Adicionalmente, el esquema de entrenamiento multi-instancia exhibió un rendimiento considerablemente inferior al de los agentes especializados, lo que restringe por el momento la utilidad práctica de un modelo único de propósito general bajo la arquitectura actual. Finalmente, es necesario señalar que los experimentos pertenecientes al eje de generalización se ejecutaron utilizando una única semilla por configuración, apartándose del protocolo de cinco semillas aplicado en el resto de la evaluación formal. Si bien los hallazgos cualitativos derivados de esta fase exploratoria se consideran metodológicamente robustos al ser verificados a nivel de mecanismo interno de decisión y no solo mediante el resultado agregado, su validación estadística formal con múltiples semillas se establece como una tarea pendiente para futuras iteraciones del trabajo.

En consecuencia, la presente investigación consolida cuatro contribuciones fundamentales al estudio del Problema de Ruteo de Vehículos Periódico mediante la aplicación de técnicas de Aprendizaje por Refuerzo. La primera de ellas radica en el desarrollo de una formulación matemática completa y plenamente funcional del PVRP estructurado como un proceso de decisión de Markov. Esta formulación integra un diseño de representación de estado holístico, un espacio de acciones discreto que incorpora de manera explícita la decisión operativa de cerrar una ruta, y un robusto mecanismo de enmascaramiento que garantiza la factibilidad local en cada iteración. En conjunto, este diseño demuestra que los enfoques basados en datos poseen la capacidad de abordar variantes del VRP que incluyen una dimensión temporal compleja, todo ello sin la necesidad de recurrir a arquitecturas de red altamente especializadas ni requerir un rediseño exhaustivo de los algoritmos de aprendizaje subyacentes. 

La segunda contribución consiste en la identificación y caracterización empírica respecto a la importancia crítica que reviste el diseño de la función de recompensa. Se comprobó que la adopción de una penalización estrictamente proporcional a la cantidad de clientes no atendidos, en contraposición a la aplicación de una penalización de magnitud estática, resultó ser un factor determinante para orientar al agente neuronal hacia la construcción de soluciones operativas completas. El episodio de vulneración de recompensa documentado durante la fase de desarrollo constituye una evidencia concreta y verificable de este principio rector, el cual resulta metodológicamente transferible a otros problemas de optimización combinatoria que presenten restricciones globales análogas. 

Como tercera contribución, se destaca la caracterización sistemática del comportamiento algorítmico del método segmentado por familia estructural del PVRP. La identificación formal de las tres familias analizadas y su posterior evaluación intensiva permiten establecer una discusión profunda sobre las capacidades y los límites del enfoque en términos de complejidad combinatoria pura, trascendiendo el análisis simplista basado únicamente en el tamaño de la red de clientes. Este nivel de categorización analítica resulta de gran utilidad para encauzar futuros trabajos de investigación que busquen posicionar sus avances dentro de un régimen operativo específico, evitando así la presentación de mejoras algorítmicas genéricas que carezcan de un alcance debidamente acotado. 

Finalmente, la cuarta contribución aporta evidencia empírica novedosa, obtenida mediante un análisis a nivel de mecanismo interno y no limitada a la simple observación del resultado final agregado, respecto a la tensión existente entre la especialización y la generalización en modelos de RL. Los experimentos diseñados en torno a la transferencia directa de políticas y al entrenamiento multi-instancia demuestran de forma cuantitativa y verificable las condiciones exactas bajo las cuales el aprendizaje se transfiere de manera genuina. Asimismo, exponen con claridad los escenarios en los que un resultado numéricamente favorable puede resultar estadísticamente engañoso, o bien, evidencian cómo la generalización forzada termina por degradar severamente el desempeño global del modelo.

Los hallazgos de esta investigación abren diversas líneas de trabajo futuras concretas y accionables a partir de la infraestructura tecnológica ya desarrollada. La proyección más directa consiste en el diseño de un esquema híbrido que combine la fase constructiva del agente con un refinamiento local, siguiendo la línea propuesta por~\cite{Chen}. Este enfoque abordaría simultáneamente los cruces de aristas mediante operadores heurísticos de tipo $2$-opt y repararía las soluciones incompletas de la familia B reinsertando a los clientes abandonados, aprovechando al agente ya entrenado como generador inicial sin requerir un reentrenamiento desde cero. Asimismo, la degradación observada en el entrenamiento multi-instancia sugiere que la arquitectura densa actual no captura eficientemente la variabilidad simultánea entre configuraciones. Por lo tanto, la incorporación de arquitecturas atencionales, tales como transformers o redes neuronales sobre grafos, podría representar de mejor manera la estructura relacional del problema, permitiendo consolidar un agente único y competitivo frente a los modelos especializados. El alcance metodológico también podría ampliarse evaluando el sistema sobre variantes del problema que incluyan ventanas de tiempo, demandas estocásticas o múltiples depósitos, extensiones que el entorno de simulación admite de forma natural gracias a su diseño modular. Por otra parte, la infactibilidad detectada en $p02$ motiva la realización de un estudio dedicado exclusivamente a la familia B que evalúe las instancias $p05$ y $p08$. Esto permitiría analizar qué configuraciones paramétricas específicas provocan el quiebre del modelo, caracterizando con precisión el umbral de factibilidad del método. Finalmente, el episodio de vulneración de recompensa documentado sugiere que el diseño de la señal de refuerzo merece una investigación independiente. La exploración de formulaciones alternativas, como recompensas intermedias basadas en la cobertura acumulada o penalizaciones dependientes de la restricción violada, podría mejorar sustancialmente el rendimiento del agente en escenarios de alta complejidad y aportar a una comprensión sistemática sobre el diseño de señales robustas en optimización combinatoria.

En definitiva, la presente investigación demuestra que el aprendizaje por refuerzo constituye una alternativa competitiva para la resolución del PVRP dentro de un régimen específico y bien caracterizado. En instancias de tamaño moderado con una estructura combinatoria simple, el agente propuesto supera a métodos clásicos reconocidos logrando tiempos de respuesta órdenes de magnitud menores. Al mismo tiempo, la evaluación exhaustiva reveló que el enfoque constructivo puro pierde competitividad al escalar la red, que existen configuraciones combinatorias donde el método no alcanza la factibilidad, y que la capacidad de generalización del modelo no puede asumirse exclusivamente a partir de un resultado numérico favorable sin antes inspeccionar el comportamiento algorítmico que lo produjo. Más allá de las cifras, el desarrollo de este trabajo sugiere que el verdadero valor de los enfoques data-driven en la optimización combinatoria no reside en su capacidad de reemplazar por completo a las metaheurísticas clásicas, sino en su enorme potencial para complementarlas. En este sentido, los sistemas híbridos que logran acoplar la construcción rápida de un agente inteligente con el refinamiento riguroso de la búsqueda local constituyen el camino más productivo para las próximas generaciones de plataformas de planificación logística. Frente a un contexto de creciente complejidad y dinamismo dentro de las cadenas de suministro globales, la contribución final de esta memoria consiste en exponer con rigor empírico y transparencia científica tanto lo que un modelo de aprendizaje puede aportar hoy al problema del ruteo periódico como las direcciones concretas hacia las cuales debe evolucionar, documentando con detalle sus aciertos operativos, sus límites estructurales y las valiosas lecciones metodológicas adquiridas durante el proceso.